\documentclass[11pt,a4paper]{article}

\usepackage[T1]{fontenc}
\usepackage[utf8]{inputenc}
\usepackage[english]{babel}
\usepackage{lmodern}
\usepackage{microtype}
\usepackage[a4paper,margin=20mm,headheight=15pt]{geometry}
\usepackage{graphicx}
\usepackage{booktabs}
\usepackage{tabularx}
\usepackage{longtable}
\usepackage{array}
\usepackage{xcolor}
\usepackage{amsmath}
\usepackage{caption}
\usepackage{float}
\usepackage{enumitem}
\usepackage{fancyhdr}
\usepackage[hidelinks]{hyperref}
\usepackage{xurl}

\graphicspath{{figures/}}
\definecolor{baselineblue}{HTML}{2563EB}
\definecolor{balancedgreen}{HTML}{16A34A}
\definecolor{maximumred}{HTML}{DC2626}
\definecolor{slate}{HTML}{334155}
\definecolor{lightgray}{HTML}{F1F5F9}

\newcommand{\dstwr}{DS-TWR}
\newcommand{\code}[1]{\texttt{\detokenize{#1}}}
\newcolumntype{Y}{>{\raggedright\arraybackslash}X}

\pagestyle{fancy}
\fancyhf{}
\lhead{\small Native DS-TWR speed-limit study}
\rhead{\small July 26, 2026}
\cfoot{\thepage}

\hypersetup{
  pdftitle={Native DS-TWR speed-limit study},
  pdfauthor={UWB ESP-IDF technical report},
  pdfsubject={Static LOS throughput, reliability, and accuracy sweep},
  pdfkeywords={UWB, DS-TWR, DW3000, ESP32-S3, ranging, latency}
}

\begin{document}
\pagenumbering{gobble}

\begin{titlepage}
  \centering
  \vspace*{17mm}
  {\Huge\bfseries Native \dstwr{} speed-limit study\par}
  \vspace{4mm}
  {\Large Throughput, reliability, and unfiltered static accuracy\par}
  \vspace{16mm}

  \begin{tabular}{rl}
    Test date: & July 26, 2026, EEST\\
    Platform: & 5 ESP32-S3 + DW3000 modules, Raspberry Pi 5\\
    Branch: & \texttt{experiment/ds-twr-native-performance}\\
    Firmware baseline: & \texttt{a97cdce-dirty}\\
    Accelerated firmware: & \texttt{27a28b0-dirty} + instrumentation patch\\
    Configuration: & M1 tag; M2, M3, M4, M5 anchors\\
    Geometry: & surveyed $3\times3$ m square, tag at the center\\
    Environment: & static, line of sight\\
  \end{tabular}

  \vfill
  \colorbox{lightgray}{%
    \begin{minipage}{0.88\textwidth}
      \vspace{3mm}
      \textbf{Key result.}
      The validated 17 ms profile delivered \textbf{58.02 complete
      positions/s}, \textbf{99.85\% frame completeness}, and
      \textbf{3.35 cm 2D RMSE}: a \textbf{3.83$\times$} speedup over the
      64 ms reference without a measurable static-accuracy penalty.
      The tested scheduling limit was 13 ms, which delivered
      \textbf{74.68 positions/s} and \textbf{3.47 cm RMSE}, but with much
      less timing margin and about 9.4 times more failed exchanges than the
      17 ms profile.
      \vspace{3mm}
    \end{minipage}}
  \vfill
  {\large Reproducible field report\par}
  \vspace{5mm}
\end{titlepage}

\pagenumbering{roman}
\tableofcontents
\clearpage
\pagenumbering{arabic}

\section{Executive summary}

This study asks how far the current native three-message \dstwr{}
implementation can be accelerated without sacrificing the static accuracy
observed in the earlier 64 ms field test. The radio exchange remains strictly
\texttt{POLL--RESP--FINAL}; no additional ranging frames, temporal filters,
moving averages, or position smoothing were introduced.

The sweep covered frame periods from 64 ms down to 13 ms. Six 60 s screening
blocks identified the useful region. The two fastest candidates were then
validated in the interleaved order
\texttt{13 ms V1--17 ms V1--13 ms V2--17 ms V2}, with two 120 s blocks per
profile. Together with the immutable 3$\times$120 s reference, the eligible
data contain 1,200.6 s of measurements, 218,617 ranges, and 54,214 complete
positions.

\begin{table}[H]
  \centering
  \caption{Main results. The 17 and 13 ms rows use pooled data from their two
  120 s validation blocks.}
  \small
  \begin{tabular}{lrrrrrr}
    \toprule
    Profile & Position rate & Speedup & Complete & RMSE & P95 & Radio success\\
    \midrule
    64 ms reference & 15.15 Hz & 1.00$\times$ & 97.34\% & 3.25 cm & 4.80 cm & n/a\\
    21 ms screening & 47.10 Hz & 3.11$\times$ & 99.89\% & 3.56 cm & 5.03 cm & 99.982\%\\
    \textcolor{balancedgreen}{17 ms balanced}
      & 58.02 Hz & 3.83$\times$ & 99.85\% & 3.35 cm & 4.65 cm & 99.966\%\\
    \textcolor{maximumred}{13 ms maximum}
      & 74.68 Hz & 4.93$\times$ & 98.71\% & 3.47 cm & 4.83 cm & 99.678\%\\
    \bottomrule
  \end{tabular}
\end{table}

\paragraph{Decision.}
The 13 ms profile is the measured throughput limit of the current
millisecond-resolution scheduler. It remains usable in this test, but its
3 ms slots have little reserve for host-task jitter. The 17 ms profile is the
recommended operating point: it preserves essentially all of the measured
accuracy, loses only 0.15\% of frames, and leaves substantially more timing
headroom. The modules were left on this 17 ms profile.

\section{System and protocol}

\subsection{Native three-frame DS-TWR}

For each anchor, the tag transmits \texttt{POLL}, the anchor returns
\texttt{RESP} using delayed transmission, and the tag sends
\texttt{FINAL}. The final frame contains the tag-side hardware timestamps;
the anchor already owns the complementary receive/transmit timestamps and
therefore calculates the range without a report or acknowledgement frame.
This is the native asymmetric double-sided two-way-ranging structure
described by Decawave and used by the Qorvo examples~\cite{aps013,qorvo-api}.

The time of flight is estimated as

\begin{equation}
  \widehat{T}_{\mathrm{prop}} =
  \frac{R_aR_b-D_aD_b}{R_a+D_a+R_b+D_b},
\end{equation}

where $R_a,R_b$ are round-trip intervals and $D_a,D_b$ are reply delays
formed from DW3000 hardware timestamps. Antenna-delay calibration and the
instantaneous carrier-integrator correction remain enabled. These are
physical timestamp corrections, not temporal filtering.

\subsection{Frame scheduler}

The tag visits anchors A2, A3, A4, and A5 sequentially. For four anchors,

\begin{equation}
  T_{\mathrm{frame}} = 4T_{\mathrm{slot}} + T_{\mathrm{gap}},
  \qquad
  f_{\mathrm{ideal}} = T_{\mathrm{frame}}^{-1}.
\end{equation}

Each complete position requires all four ranges associated with the same
sequence round. The 64--21 ms profiles retain 2 ms response and final delays.
The 17 and 13 ms profiles use 1 ms for each delayed transmission. Shortening
these delays changes only the scheduled turnaround time, not the messages or
the DS-TWR estimator.

\begin{table}[H]
  \centering
  \caption{Timing profiles tested. All values except the ideal rate are in
  milliseconds.}
  \small
  \begin{tabular}{rrrrrrr}
    \toprule
    Frame & Slot & Gap & RX timeout & RESP delay & FINAL delay & Ideal Hz\\
    \midrule
    64 & 15 & 4 & 8 & 2 & 2 & 15.63\\
    33 & 8 & 1 & 5 & 2 & 2 & 30.30\\
    29 & 7 & 1 & 5 & 2 & 2 & 34.48\\
    25 & 6 & 1 & 5 & 2 & 2 & 40.00\\
    21 & 5 & 1 & 5 & 2 & 2 & 47.62\\
    17 & 4 & 1 & 4 & 1 & 1 & 58.82\\
    13 & 3 & 1 & 3 & 1 & 1 & 76.92\\
    \bottomrule
  \end{tabular}
\end{table}

\section{Experimental method}

\subsection{Geometry and reference data}

The anchors form a nominal $3\times3$ m square and the tag is at
$(1.5,1.5)$ m. Physical placement uncertainty is approximately
$\pm2$ mm. The true center-to-corner range is
$\sqrt{1.5^2+1.5^2}=2.12132$ m. All equipment remained static and in line of
sight throughout the sweep.

The 64 ms reference is the three-block DS-TWR data set from the preceding
protocol comparison. It contains 360 s, 22,256 ranges, and 5,455 complete
positions. Reusing it avoids redefining the baseline after seeing accelerated
results.

\subsection{Collection sequence}

Every mode change was applied to all five modules in parallel and followed by
a reboot and a warm-up interval. The eligible accelerated sequence was:

\begin{enumerate}[leftmargin=8mm]
  \item 60 s screening at 33, 29, 25, 21, 17, and 13 ms;
  \item 120 s at 13 ms, then 120 s at 17 ms;
  \item a second 120 s block at 13 ms, then a second at 17 ms.
\end{enumerate}

The interleaved validation order limits the chance that temperature or slow
environmental drift favors one final candidate.

\subsection{Instrumentation controls}

The first exploratory 33 ms run revealed a deterministic instrumentation
artifact: extended DW3000 diagnostic-register reads and verbose log output
executed periodically in the anchor's real-time path. They caused the next
poll to be missed in an almost exact five-sequence pattern. The firmware was
changed once so that continuous positioning retains one compact result record
per successful exchange while extended diagnostics remain available in the
dedicated distance-test and survey modes. No protocol frame or ranging
equation changed. Repeating the identical 33 ms profile improved frame
completeness from 19.41\% to 98.62\%.

At 21 ms, the original collector reached a separate observation limit because
it read only the latest distance per anchor from an HTTP snapshot. The radio
still reported 99.982\% successful exchanges. Collection was changed to an
ordered log cursor, then the same profile was repeated. The valid repeat
captured 99.89\% complete frames.

Both exploratory runs are retained verbatim and marked
\texttt{comparison\_eligible=false}. Their exclusion was based on known
instrumentation/collection behavior, not on range or position error.

\subsection{Metrics}

Ranges are grouped into rounds using the anchor index and the 16-bit sequence
number. A frame is complete only if all four anchors are present. Positions
are reconstructed by the same unfiltered linear trilateration used for the
reference report. The reported error is Euclidean distance from
$(1.5,1.5)$ m.

Radio success, timeouts, maximum exchange duration, and slot overruns come
from one-second tag summaries and failure records. The historical baseline
capture predates this timing-log collector, so its radio-success entry is
shown as unavailable rather than inferred.

\section{Results}

\subsection{Throughput and completeness}

\begin{figure}[H]
  \centering
  \includegraphics[width=\textwidth]{01_rate_completeness.pdf}
  \caption{Effective complete-position rate versus the theoretical scheduler
  rate, and frame completeness. Final 17 and 13 ms values use validation
  blocks; intermediate profiles are screening results.}
\end{figure}

The 17 ms profile reaches 98.6\% of its theoretical rate and produces
3.83 times the reference update rate. The 13 ms profile reaches 97.1\% of
theoretical and produces 4.93 times the reference update rate. Even at the
tested limit, its 98.71\% completeness remains above the 97.34\% historical
reference value.

\subsection{Static accuracy}

\begin{figure}[H]
  \centering
  \includegraphics[width=0.92\textwidth]{02_accuracy.pdf}
  \caption{Unfiltered position RMSE, median, and P95 across the timing sweep.}
\end{figure}

There is no monotonic loss of accuracy with update rate. Relative to the
reference, the validated 17 ms profile changes RMSE by only +0.99 mm and
improves P95 by 1.45 mm. The 13 ms profile changes RMSE by +2.17 mm and P95
by +0.33 mm. These differences are comparable to placement uncertainty and
between-block variation and should not be interpreted as a causal accuracy
change.

\begin{figure}[H]
  \centering
  \includegraphics[width=0.90\textwidth]{04_position_error_cdf.pdf}
  \caption{Empirical CDF of individual unfiltered position errors for the
  reference and the two validated accelerated profiles.}
\end{figure}

The distributions substantially overlap. The maximum errors were 6.79 cm at
64 ms, 7.70 cm at 17 ms, and 7.20 cm at 13 ms; no large outlier comparable to
the previously excluded pathological FlexTDOA event appears.

\subsection{Reliability and timing margin}

\begin{figure}[H]
  \centering
  \includegraphics[width=0.94\textwidth]{03_reliability.pdf}
  \caption{Failed exchanges and slot overruns. The final two profiles are
  derived from 240 s each; intermediate profiles are 60 s screening blocks.}
\end{figure}

\begin{table}[H]
  \centering
  \caption{Reliability of the final candidates, pooled over two 120 s blocks.}
  \small
  \begin{tabular}{lrrrrr}
    \toprule
    Profile & Attempts & Failures & Success & Failures/1,000 & Overruns/min\\
    \midrule
    17 ms & 55,640 & 19 & 99.966\% & 0.341 & 4.50\\
    13 ms & 72,604 & 234 & 99.678\% & 3.223 & 61.96\\
    \bottomrule
  \end{tabular}
\end{table}

The faster profile gains 16.66 complete positions/s, but failed exchanges are
about 9.4 times more frequent and slot overruns about 13.8 times more
frequent. This is the clearest evidence that 13 ms is a boundary profile
rather than a comfortable operating point.

\begin{figure}[H]
  \centering
  \includegraphics[width=0.94\textwidth]{07_timing_margin.pdf}
  \caption{Distribution of the maximum exchange duration in failure-free
  one-second windows, compared with the configured slot budget.}
\end{figure}

At 17 ms, the median and P95 one-second maxima are 2.63 and 2.73 ms inside a
4 ms slot; even P99 is 3.48 ms. At 13 ms, the median is similar at 2.62 ms,
but P95 rises to 3.46 ms and P99 to 4.74 ms, both beyond the 3 ms slot.
Occasional scheduling/logging jitter therefore creates frequent overruns at
the limit. A shorter nominal frame with the current integer-millisecond
configuration would overrun by construction and was not treated as a valid
profile.

\subsection{Range bias and repeatability}

\begin{figure}[H]
  \centering
  \includegraphics[width=0.94\textwidth]{05_anchor_bias.pdf}
  \caption{Per-anchor mean range bias over the sweep. The persistent
  module-dependent calibration offsets dominate any timing-profile effect.}
\end{figure}

Anchor biases remain approximately +2.4 to +3.2 cm for A2, +7.7 to +8.0 cm
for A3, +2.5 to +3.0 cm for A4, and $-4.5$ to $-4.1$ cm for A5. Their
stability explains why accelerated position accuracy stays close to the
reference: response-delay reduction does not introduce a new systematic
range shift in this data.

\begin{figure}[H]
  \centering
  \includegraphics[width=\textwidth]{06_validation_repeatability.pdf}
  \caption{The two interleaved 120 s validation blocks for each final
  profile.}
\end{figure}

The validation blocks reproduce both rate and accuracy. The 13 ms position
rates are 74.60 and 74.76 Hz; the 17 ms rates are 58.01 and 58.03 Hz. RMSE
varies by 1.10 mm between the two 13 ms blocks and by 0.27 mm between the two
17 ms blocks.

\section{Interpretation and recommendation}

\subsection{What limits the current implementation}

The initial 64 ms configuration was dominated by deliberate idle time, not
UWB airtime. Removing nonessential diagnostics from the real-time path and
reducing scheduled turnaround delays exposes the actual implementation
envelope:

\begin{itemize}[leftmargin=7mm]
  \item with 2+2 ms delayed transmissions, 21 ms is a strong operating point
  at 47.10 positions/s;
  \item with 1+1 ms delayed transmissions, 17 ms is robust and 13 ms is the
  current scheduling boundary;
  \item the 13 ms boundary is set by the 3 ms host slot, integer-millisecond
  runtime configuration, and occasional ESP32 task/log jitter---not by a
  demonstrated degradation of the DS-TWR distance estimator.
\end{itemize}

Moving below 13 ms cleanly would require sub-millisecond runtime timing,
tighter real-time scheduling, or a design that pipelines work differently.
That would be a separate firmware optimization study. Simply configuring a
2 ms slot today would make most healthy exchanges overrun their nominal slot
and would not constitute a valid 9 ms frame.

\subsection{Recommended profiles}

\paragraph{Default: 17 ms balanced profile.}
Use slot 4 ms, round gap 1 ms, RX timeout 4 ms, response delay 1 ms, final
delay 1 ms, and auto-RX delay 500 UUS. It gives 58.02 complete positions/s,
99.85\% completeness, and accuracy statistically indistinguishable from the
64 ms reference, while retaining useful scheduling reserve.

\paragraph{Maximum-throughput option: 13 ms.}
Use slot 3 ms, gap 1 ms, timeout 3 ms, and 1+1 ms delayed transmissions when
74.68 positions/s is worth the reduced timing margin. The test remained
accurate and 98.71\% complete, but this setting should be revalidated under
network load, motion, and longer operation before deployment.

\paragraph{Conservative option: 21 ms.}
If retaining the original 2+2 ms UWB reply delays is desirable, the 21 ms
profile supplies 47.10 positions/s and 99.89\% frame completeness.

\section{Limits of inference and next tests}

This is a rigorous speed sweep at one favorable static point, not a complete
localization qualification. The conclusions apply to the tested modules,
firmware, geometry, channel, and line-of-sight environment. They do not yet
establish performance:

\begin{itemize}[leftmargin=7mm]
  \item near edges or outside the anchor square;
  \item while the tag is moving or rotating;
  \item under human blockage, multipath, or NLOS;
  \item with different anchor/tag antenna orientations;
  \item under heavy simultaneous Wi-Fi/dashboard load;
  \item over hours of thermal drift or with multiple tags.
\end{itemize}

The next study should use a spatial grid with independently measured ground
truth, repeat 64/21/17/13 ms in randomized order at every point, and add
controlled dynamic trajectories. The 17 ms profile is the appropriate
starting point for those tests.

\clearpage
\appendix

\section{Complete profile table}

\begin{table}[H]
  \centering
  \caption{Aggregate comparison. Intermediate profiles are screening blocks;
  17 and 13 ms use validation data.}
  \scriptsize
  \begin{tabular}{rrrrrrrr}
    \toprule
    Frame & Positions/s & Complete & RMSE & Median & P95 & Success & Fail/1,000\\
    \midrule
    64 & 15.15 & 97.34\% & 3.25 cm & 3.04 cm & 4.80 cm & n/a & n/a\\
    33 & 29.67 & 98.62\% & 3.41 cm & 3.29 cm & 4.92 cm & 99.916\% & 0.835\\
    29 & 33.46 & 97.81\% & 3.43 cm & 3.28 cm & 4.92 cm & 99.976\% & 0.242\\
    25 & 38.61 & 97.44\% & 3.53 cm & 3.43 cm & 5.03 cm & 99.989\% & 0.106\\
    21 & 47.10 & 99.89\% & 3.56 cm & 3.42 cm & 5.03 cm & 99.982\% & 0.177\\
    17 & 58.02 & 99.85\% & 3.35 cm & 3.24 cm & 4.65 cm & 99.966\% & 0.341\\
    13 & 74.68 & 98.71\% & 3.47 cm & 3.36 cm & 4.83 cm & 99.678\% & 3.223\\
    \bottomrule
  \end{tabular}
\end{table}

\section{Files and reproduction}

\small
\begin{tabularx}{\textwidth}{p{57mm}Y}
  \toprule
  File/directory & Contents\\
  \midrule
  \path{data/*.jsonl.xz} & lossless archives of raw range/status events and
  timing logs\\
  \path{data/*.metadata.json} & timing configuration, module state, and
  collection metadata\\
  \path{raw_data_manifest.csv} & byte counts and SHA-256 checksums for raw and
  archived files\\
  \path{block_metrics.csv} & metrics for every baseline, screening,
  validation, and exploratory block\\
  \path{profile_metrics.csv} & phase-aware aggregate metrics\\
  \path{positions.csv} & every reconstructed complete position\\
  \path{analysis_summary.json} & machine-readable analysis output\\
  \path{figures/} & all plots in vector PDF and PNG\\
  \bottomrule
\end{tabularx}

\paragraph{Reproduction}
\begin{verbatim}
cd /home/pi/Documents/UWB
xz --decompress --keep reports/ds_twr_speed_limit_20260726/data/*.jsonl.xz
python3 tools/uwb_ds_speed_analyze.py
cd reports/ds_twr_speed_limit_20260726
latexmk -pdf -interaction=nonstopmode -halt-on-error report.tex
\end{verbatim}

The analysis never modifies raw event files. The original 64 ms blocks live
under \path{reports/uwb_protocol_comparison_20260726/data}; the path is
recorded in \path{analysis_summary.json}.

\begin{thebibliography}{9}

\bibitem{qorvo-api}
Qorvo, \emph{DW3xxx Device Driver API Guide}, double-sided two-way-ranging
examples, 2025.
\url{https://forum.qorvo.com/uploads/short-url/xD3TlXKvkujjdUaJWXv2b4E7GQN.pdf}.

\bibitem{aps013}
Decawave, \emph{APS013: DW1000 and Two-Way Ranging}, rev. 2.3,
Poll--Response--Final and the asymmetric DS-TWR formula.
\url{https://forum.qorvo.com/uploads/short-url/x34DrF7EW5fQP9wY3aNESqPKz8z.pdf}.

\bibitem{qorvo-product}
Qorvo, \emph{DW3110 Ultra-Wideband Transceiver}, product documentation.
\url{https://www.qorvo.com/products/p/DW3110}.

\bibitem{firmware}
Local repository \emph{UWB Module ESP-IDF}, branch
\texttt{experiment/ds-twr-native-performance}, reference commit
\texttt{27a28b03f1dd627cdcbbf327fb94127bb3fe371b}, working state from
July 26, 2026.

\end{thebibliography}

\end{document}
